\documentclass{article}

\usepackage{amsmath,amssymb}
\usepackage{xurl}
\usepackage[hidelinks]{hyperref}
\setlength{\emergencystretch}{2em}

% Shared commands for notes in docs/paper-gaps/.

\DeclareUnicodeCharacter{2080}{\ifmmode{}_{0}\else\textsubscript{0}\fi}
\DeclareUnicodeCharacter{2081}{\ifmmode{}_{1}\else\textsubscript{1}\fi}
\DeclareUnicodeCharacter{2082}{\ifmmode{}_{2}\else\textsubscript{2}\fi}
\DeclareUnicodeCharacter{2099}{\ifmmode{}_{n}\else\textsubscript{n}\fi}

\newcommand{\Lean}{\textsc{Lean}}
\ExplSyntaxOn
\NewDocumentCommand{\leanid}{m}{
  \ifmmode
    \textsf{\detokenize{#1}}
  \else
    \group_begin:
    \urlstyle{sf}
    \tl_set:Nn \l_tmpa_tl {#1}
    \tl_replace_all:Nnn \l_tmpa_tl {\_} {_}
    \exp_args:NV \path \l_tmpa_tl
    \group_end:
  \fi
}
\ExplSyntaxOff
\providecommand{\tr}{\operatorname{tr}}
\newcommand{\tnleanIssueBase}{https://github.com/LionSR/TNLean/issues/}
\newcommand{\tnleanPrBase}{https://github.com/LionSR/TNLean/pull/}
\newcommand{\tnleanDocBase}{https://sirui-lu.com/TNLean/docs/}
\newcommand{\ghissue}[1]{\href{\tnleanIssueBase#1}{GitHub issue \##1}}
\newcommand{\ghpr}[1]{\href{\tnleanPrBase#1}{PR \##1}}
\newcommand{\doclink}[1]{\href{\tnleanDocBase#1.html}{\path{#1}}}

% Dirac notation and the identity operator, as in blueprint/src/macros/common.tex.
% \providecommand keeps note-local definitions authoritative.
\usepackage{dsfont}
\providecommand{\ket}[1]{|#1\rangle}
\providecommand{\bra}[1]{\langle #1|}
\providecommand{\braket}[2]{\langle #1 | #2 \rangle}
\providecommand{\ketbra}[2]{|#1\rangle\!\langle #2|}
\providecommand{\Id}{\mathds{1}}
\providecommand{\one}{\mathds{1}}
\providecommand{\C}{\mathbb{C}}
\providecommand{\MN}[1]{M_{#1}(\C)}
\providecommand{\MD}{M_D(\C)}

% External referencing for paper sources.  The build helper writes sanitized
% auxiliary files under build/paper-aux/ and excludes duplicated source labels.
\usepackage{xr}

% Import a generated paper auxiliary file with a caller-supplied label prefix.
% The candidate paths support builds from docs/paper-gaps/, the repository root,
% and build/paper-aux/ itself.
\newcommand{\importpaperaux}[2]{%
  \IfFileExists{../../build/paper-aux/#2.aux}{%
    \externaldocument[#1]{../../build/paper-aux/#2}%
  }{%
    \IfFileExists{build/paper-aux/#2.aux}{%
      \externaldocument[#1]{build/paper-aux/#2}%
    }{%
      \IfFileExists{#2.aux}{%
        \externaldocument[#1]{#2}%
      }{}%
    }%
  }%
}

% General external reference macros
\newcommand{\paperref}[2]{\ref{#1-#2}}
\newcommand{\papereqref}[2]{(\ref{#1-#2})}

% CPSV16 source (1606.00608 / MPDO-22-12-17-2.tex) external document and shortcuts
\importpaperaux{cpsv16-}{1606.00608/MPDO-22-12-17-2-tnlean}
\newcommand{\cpsvref}[1]{\paperref{cpsv16}{#1}}
\newcommand{\cpsveqref}[1]{\papereqref{cpsv16}{#1}}

% Verdict marker: \gapnote{<kind>}{<status>}. Kind is the policy
% classification (clarification, local-correction, scope-restriction,
% unfaithful, false-source, open-gap); status is open, wip (elimination
% underway), resolved, or historical. Machine-read by the site index;
% prints a verdict line.
\newcommand{\gapnote}[2]{\par\noindent\textbf{Verdict:} #1 (#2).\par}


\title{Top-index scope of the spectral n-positivity criterion (Wolf Prop.~3.2)}
\author{}
\date{2026-06-24}

\begin{document}
\maketitle

\section{Source statement}

Proposition~3.2 of Wolf's \emph{Quantum Channels \& Operations}
(\path{Notes/WolfNoteTexSource/ch03_positive_not_completely.tex}, lines
153--179) bounds the expectation of a Hermitian operator
$\tau\in\mathcal{M}_{D'}\otimes\mathcal{M}_{D}$ in vectors of fixed Schmidt
rank. Writing $\tau=\sum_i\nu_i|\phi_i\rangle\langle\phi_i|$ with normalized
eigenvectors, $\rho_i=\tr_2|\phi_i\rangle\langle\phi_i|$ the reduced density on
the first factor $\mathbb{C}^{D'}$, $\nu_0$ the smallest positive eigenvalue and
$\nu$ the largest one, the infimum over normalized vectors $\psi$ of Schmidt
rank $n$ obeys
\begin{align}
  \inf_{\psi}\langle\psi|\tau|\psi\rangle
  &\ge \nu_0 + \sum_{i:\nu_i\le 0}(\nu_i-\nu_0)\,\|\rho_i\|_{(n)},
  \tag{3.7}
\end{align}
and, when a single eigenvalue $\nu_-$ is non-positive,
\begin{align}
  \inf_{\psi}\langle\psi|\tau|\psi\rangle
  &\le \nu + (\nu_- - \nu)\,\|\rho_-\|_{(n)}.
  \tag{3.8}
\end{align}
Each $\rho_i$ acts on the first factor $\mathbb{C}^{D'}$, so its Ky-Fan
$n$-norm $\|\rho_i\|_{(n)}$ is meaningful for $1\le n\le D'$. Proposition~3.1 of
the same chapter restricts the positivity index to $1\le n\le D$, the dimension
of the second factor (same file, lines 90--92), and that is the range the source
asserts. The value $n=D'$ therefore lies inside the source's range exactly when
$D'\le D$; when $D'>D$ it lies beyond, and the formal endpoint below is an
extension of the source rather than an instance of it. What matters for coverage
is that the formal statements together hold at every $n\ge1$, hence at every $n$
the source speaks about.

\section{Formalized statement}

For Schmidt ranks $1\le n<D'$ the formal theorems\footnote{Formal declarations:
    \leanid{Matrix.IsHermitian.spectral_lower_bound} and
    \leanid{Matrix.IsHermitian.exists_le_spectral_upper_bound} in
    \doclink{TNLean/Channel/NPositivitySpectralCriterion}.}
establish these bounds in the per-vector form: every normalized vector $\psi$ of
Schmidt rank at most $n$ satisfies the lower bound~(3.7), and some such $\psi$
realizes the upper bound~(3.8).

The source's infimum form of the same two bounds is stated for the infimum of
the expectations $\bra{\psi}\tau\ket{\psi}$ over the normalized vectors of
Schmidt rank at most $n$.\footnote{Formal declarations:
    \leanid{Matrix.IsHermitian.le_sInf_schmidtRankLEExpectations}~(3.7),
    \leanid{Matrix.IsHermitian.sInf_schmidtRankLEExpectations_le}~(3.8), and
    \leanid{Matrix.IsHermitian.sInf_schmidtRankLEExpectations_top} for the top
    index, in \doclink{TNLean/Channel/NPositivitySpectralCriterion}.} The
reading of the source's phrase ``vectors of Schmidt rank $n$'' is discussed in
\path{docs/paper-gaps/wolf_prop_3_2_schmidt_rank_reading.tex}.

\section{Status}

The top index $n=D'$ is resolved by a route that bypasses the $n<D'$ restriction
of the maximal-overlap and Ky-Fan principles. At $n=D'$ the Schmidt-rank
constraint is vacuous: every normalized vector qualifies and
$\|\rho_i\|_{(D')}=\tr\rho_i=1$. Both bounds therefore constrain the same quantity,
the infimum of $\langle\psi|\tau|\psi\rangle$ over all normalized vectors:
\begin{itemize}
  \item The Rayleigh estimate gives $\nu_{\min}\le\langle\psi|\tau|\psi\rangle$
  for every normalized $\psi$, so $\inf_\psi\langle\psi|\tau|\psi\rangle\ge
  \nu_{\min}$. The right-hand side of~(3.7) reads
  $\nu_0+\sum_{i:\nu_i\le0}(\nu_i-\nu_0)$ there, which is at most $\nu_{\min}$
  and in general strictly below it: for the spectrum $(-2,-1,3)$ it is $-6$
  while $\nu_{\min}=-2$. So~(3.7) at the top index follows from the Rayleigh
  estimate, but is not the same inequality.
  \item The upper bound~(3.8) on the infimum becomes the existence of a normalized
  vector with $\langle\psi|\tau|\psi\rangle=\nu_{\min}$, namely a least-eigenvalue
  eigenvector, so $\inf_\psi\langle\psi|\tau|\psi\rangle\le\nu_{\min}$; and under
  the hypothesis of~(3.8) that $\nu_-$ is the only non-positive eigenvalue, the
  right-hand side $\nu+(\nu_--\nu)\|\rho_-\|_{(D')}$ is exactly $\nu_-=\nu_{\min}$.
\end{itemize}
Together they give $\inf_\psi\langle\psi|\tau|\psi\rangle=\nu_{\min}$. Both facts
follow directly from the Rayleigh expansion $\langle\psi|\tau|\psi\rangle=\sum_i
\nu_i|\langle\phi_i|\psi\rangle|^2$ and Parseval's identity $\sum_i
|\langle\phi_i|\psi\rangle|^2=1$, comparing each $\nu_i$ against $\nu_{\min}$ and
evaluating at the least-eigenvalue eigenvector; the Ky-Fan machinery of the $n<D'$
proof is not invoked.

The top-index endpoint is recorded as
\leanid{Matrix.IsHermitian.spectral_lower_bound_top} and
\leanid{Matrix.IsHermitian.exists_spectral_lower_bound_top}, which together
state $\min_\psi\langle\psi|\tau|\psi\rangle=\nu_{\min}$. Inequality~(3.7)
itself at the top index, in the infimum form, is
\leanid{Matrix.IsHermitian.le_sInf_schmidtRankLEExpectations_top}, and~(3.8) is
the corresponding half of
\leanid{Matrix.IsHermitian.sInf_schmidtRankLEExpectations_top}.

\section{Classification}

\emph{Resolved.} Bounds~(3.7) and~(3.8) are now formalized at every index
$n\ge1$, hence over the source's range $1\le n\le D$ whatever the two dimensions
are, and beyond it when $D'>D$: the per-vector bounds
\leanid{Matrix.IsHermitian.spectral_lower_bound} and
\leanid{Matrix.IsHermitian.exists_le_spectral_upper_bound} for $1\le n<D'$, and
their endpoints
\leanid{Matrix.IsHermitian.spectral_lower_bound_top} and
\leanid{Matrix.IsHermitian.exists_spectral_lower_bound_top} for $n\ge D'$, by
the route detailed in the Status above. The infimum form of the source's
statement is available over the same range,
\leanid{Matrix.IsHermitian.le_sInf_schmidtRankLEExpectations_top} carrying~(3.7)
at the top indices. No hypothesis foreign to the source is used.

\bibliographystyle{alpha}
\bibliography{references}

\end{document}
